\documentclass[11pt,a4paper]{article}

\usepackage[margin=22mm,top=19mm,bottom=19mm,headheight=14pt,headsep=6mm,footskip=10mm]{geometry}
\usepackage{amsmath,amssymb,mathtools}
\usepackage{enumitem}
\usepackage{fancyhdr}
\usepackage{xcolor}
\usepackage{lastpage}
\usepackage[T1]{fontenc}
\usepackage{lmodern}

\definecolor{examblue}{HTML}{23395B}
\definecolor{rulegray}{HTML}{CBD2DA}

\pagestyle{fancy}
\fancyhf{}
% \lhead{\small Linear Algebra -- Draft Final Examination}
% \rhead{\small Post-midterm emphasis}
% \cfoot{\small Page \thepage\ of \pageref{LastPage}}
% \renewcommand{\headrulewidth}{0.4pt}
% \renewcommand{\footrulewidth}{0.4pt}

\setlength{\parindent}{0pt}
\setlength{\parskip}{4pt}
\setlist[enumerate]{leftmargin=*,itemsep=3pt,topsep=3pt}
\setlist[enumerate,2]{label=\Alph*.,leftmargin=9mm,itemsep=1pt,topsep=1pt}

\newcommand{\parttitle}[2]{%
  \vspace{1mm}
  {\color{examblue}\Large\bfseries #1}\quad
  {\color{examblue}\large #2}\par\vspace{3mm}
}

\newcommand{\answerbox}[1]{%
  \par\vspace{2mm}
  {\color{rulegray}\noindent\rule{\linewidth}{0.4pt}\par}
  \vspace{#1}
  {\color{rulegray}\noindent\rule{\linewidth}{0.4pt}\par}
}

\newcommand{\markpts}[1]{\hfill\textbf{(#1 marks)}}

\begin{document}

\begin{center}
  {\color{examblue}\LARGE\bfseries LINEAR ALGEBRA}\par
  \vspace{1mm}
  {\color{examblue}\Large\bfseries DRAFT FINAL EXAMINATION}
\end{center}

\vspace{3mm}
\begin{tabular*}{\textwidth}{@{\extracolsep{\fill}}ll}
Name: \rule{68mm}{0.4pt} & Student ID: \rule{48mm}{0.4pt}\\[3mm]
Time allowed: \rule{25mm}{0.4pt} & Total: 100 marks
\end{tabular*}

\vspace{4mm}
\textbf{Instructions.} Show sufficient work in Part III. Unless instructed otherwise,
give exact answers. The examination emphasizes material taught after the midterm.

\parttitle{Part I.}{Fill in the blanks. (20 marks; 2 marks each)}

\begin{enumerate}[label=\arabic*.]
  \item Consider the basis
        $\left\{\begin{bmatrix}1\\2\end{bmatrix},
        \begin{bmatrix}2\\-1\end{bmatrix}\right\}$ for $\mathbb R^{2}$.
        The coordinate vector of $\begin{bmatrix}5\\0\end{bmatrix}$ with respect to this basis is
        $\underline{\hspace{28mm}}$.
  \item If $A$ is invertible, then $\det(A^{-1})=\underline{\hspace{35mm}}$.
  \item If $A=\begin{bmatrix}2&-1\\4&3\end{bmatrix}$, then
        $A^{T}=\underline{\hspace{35mm}}$.
  \item If $A$ has 7 columns and $\operatorname{rank}(A)=4$, then
        the dimension of the null space of $(A)$ is $\underline{\hspace{22mm}}$.
  % \item The eigenspace corresponding to an eigenvalue $\lambda$ is
  %       the null space of $\underline{\hspace{30mm}}$.
  % \item The eigenvalues of $\begin{bmatrix}3&1\\0&-2\end{bmatrix}$ are
  %       $\underline{\hspace{35mm}}$.
  % \item $\det\!\begin{bmatrix}4&-2\\1&3\end{bmatrix}
  %       =\underline{\hspace{24mm}}$.
  % \item $\begin{bmatrix}-2\\3\\6\end{bmatrix}=$ 
  %       $\underline{\hspace{24mm}}$.
  % \item The distance between $\begin{bmatrix}1\\-2\\3\end{bmatrix}$ and
  %       $\begin{bmatrix}-2\\2\\3\end{bmatrix}$ is $\underline{\hspace{24mm}}$.
  % \item If $A=\begin{bmatrix}2&0&0\\0&-3&0\\0&0&4\end{bmatrix}$, then
  %       $A^{2}=\underline{\hspace{42mm}}$.
\end{enumerate}

\newpage
\parttitle{Part II.}{Select the only correct statement. (20 marks; 2 marks each)}

\begingroup\small
\begin{enumerate}[label=\arabic*.,start=11]
  \item For a square matrix $A$, which condition is equivalent to $A$ being invertible?
    \begin{enumerate}
      \item $A\mathbf x=\mathbf 0$ has a nontrivial solution.
      \item $A\mathbf x=\mathbf b$ has a unique solution for every $\mathbf b$.
      \item The columns of $A$ are linearly dependent.
      \item $\det(A)=0$.
    \end{enumerate}

  \item If $B$ is obtained from $A$ by interchanging two rows, then
    \begin{enumerate}
      \item $\det(B)=\det(A)$.
      \item $\det(B)=-\det(A)$.
      \item $\det(B)=2\det(A)$.
      \item $\det(B)=0$ for every $A$.
    \end{enumerate}

  \item Which set is a subspace of $\mathbb R^{3}$?
    \begin{enumerate}
      \item $\{(x,y,z):x-2y+z=0\}$
      \item $\{(x,y,z):x-2y+z=1\}$
      \item $\{(x,y,z):x,y,z>0\}$
      \item $\{(x,y,z):x+y+z\geq 0\}$
    \end{enumerate}

  \item Pivot columns 1 and 3 occur in $\operatorname{rref}(A)$. A basis for
        $\operatorname{Col}(A)$ is obtained from
    \begin{enumerate}
      \item columns 1 and 3 of the original matrix $A$.
      \item columns 1 and 3 of $\operatorname{rref}(A)$.
      \item rows 1 and 3 of $A$.
      \item all nonzero rows of $\operatorname{rref}(A)$.
    \end{enumerate}

  \item If $A$ has 5 columns and $\operatorname{rank}(A)=3$, then
        $\dim\operatorname{Null}(A)$ equals
    \begin{enumerate}
      \item 1 \item 2 \item 3 \item 5
    \end{enumerate}
\end{enumerate}
\endgroup

\newpage
\parttitle{Part II.}{Continued}

\begin{enumerate}[label=\arabic*.,start=16]
  \item Eigenvectors corresponding to distinct eigenvalues of a matrix are always
    \begin{enumerate}
      \item orthogonal. \item linearly independent.
      \item equal in length. \item elements of the null space of the matrix.
    \end{enumerate}

  \item Every real symmetric matrix is
    \begin{enumerate}
      \item singular. \item upper triangular.
      \item orthogonally diagonalizable. \item guaranteed to have distinct eigenvalues.
    \end{enumerate}

  \item If $Q$ is an orthogonal matrix, then
    \begin{enumerate}
      \item $Q^{-1}=Q^{T}$. \item $Q^{T}Q=0$.
      \item $Q^{-1}=-Q$. \item the columns of $Q$ are linearly dependent.
    \end{enumerate}

  \item If $\widehat{\mathbf y}=\operatorname{proj}_{W}\mathbf y$, then
        $\mathbf y-\widehat{\mathbf y}$ is
    \begin{enumerate}
      \item in $W$. \item orthogonal to $W$.
      \item equal to $\mathbf 0$ for every $\mathbf y$. \item parallel to every vector in $W$.
    \end{enumerate}

  \item For a least-squares solution $\widehat{\mathbf x}$, the residual
        $\mathbf b-A\widehat{\mathbf x}$ is orthogonal to
    \begin{enumerate}
      \item $\operatorname{Null}(A)$. \item $\operatorname{Col}(A)$.
      \item $\operatorname{Row}(A)$. \item $\operatorname{Null}(A^{T})$.
    \end{enumerate}
\end{enumerate}

\vspace{4mm}
{\color{examblue}\bfseries End of Part II. Continue to Part III and show all work.}

\newpage
\parttitle{Part III.}{Show all your work in the space provided. (60 marks)}

\textbf{21.} Use Gaussian elimination to solve the system
\[
\begin{aligned}
x+y+z&=2,\\
2x-y+z&=-1,\\
x+2y-z&=6.
\end{aligned}
\]
State whether the solution is unique and justify your conclusion. \markpts{8}

\textbf{Solution:}
\answerbox{100mm}

\newpage
\textbf{22.} Let
\[
A=\begin{bmatrix}
1&0&1&2\\
0&1&1&-1\\
1&1&2&1
\end{bmatrix}.
\]
\begin{enumerate}[label=(\alph*),leftmargin=9mm]
  \item Find $\operatorname{rref}(A)$ and identify the pivot columns. \markpts{4}
  \item Find a basis for $\operatorname{Null}(A)$. \markpts{4}
  \item Find a basis for $\operatorname{Col}(A)$. \markpts{3}
  \item Verify the rank-nullity theorem for $A$. \markpts{3}
\end{enumerate}

\textbf{Solution:}
\answerbox{130mm}

\newpage
\textbf{23.} Consider
\[
B=\begin{bmatrix}4&1\\2&3\end{bmatrix}.
\]
\begin{enumerate}[label=(\alph*),leftmargin=9mm]
  \item Find the eigenvalues of $B$ and a basis for each eigenspace. \markpts{6}
  \item Diagonalize $B$ by finding matrices $P$ and $D$ such that $B=PDP^{-1}$. \markpts{4}
  \item Use the diagonalization to compute $B^{4}$. \markpts{2}
\end{enumerate}

\textbf{Solution:}
\answerbox{120mm}

\newpage
\textbf{24.} Let
\[
\mathbf v_{1}=\begin{bmatrix}1\\1\\0\end{bmatrix},\qquad
\mathbf v_{2}=\begin{bmatrix}1\\0\\1\end{bmatrix},\qquad
W=\operatorname{span}\{\mathbf v_{1},\mathbf v_{2}\},\qquad
\mathbf b=\begin{bmatrix}2\\1\\3\end{bmatrix}.
\]
\begin{enumerate}[label=(\alph*),leftmargin=9mm]
  \item Apply the Gram-Schmidt process to obtain an orthonormal basis for $W$. \markpts{6}
  \item Find $\operatorname{proj}_{W}\mathbf b$. \markpts{4}
  \item Verify that $\mathbf b-\operatorname{proj}_{W}\mathbf b$ is orthogonal to $W$. \markpts{2}
\end{enumerate}

\textbf{Solution:}
\answerbox{120mm}

\newpage
\textbf{25.} A straight line $y=\beta_{0}+\beta_{1}t$ is to be fitted, in the
least-squares sense, to the four data points
\[
(0,1),\qquad(1,2),\qquad(2,2),\qquad(3,4).
\]
\begin{enumerate}[label=(\alph*),leftmargin=9mm]
  \item Write the system $A\boldsymbol\beta\approx\mathbf y$ and form the normal equations. \markpts{4}
  \item Solve the normal equations to find the least-squares line. \markpts{5}
  \item Use the fitted line to predict $y$ when $t=4$. \markpts{2}
  \item Compute the residual vector and verify that it is orthogonal to
        $\operatorname{Col}(A)$. \markpts{3}
\end{enumerate}

\textbf{Solution:}
\answerbox{50mm}

\newpage
{\color{examblue}\large\bfseries Question 25 continued / Additional workspace}

\vspace{4mm}
\answerbox{200mm}

\vfill
\begin{center}
  {\color{examblue}\Large\bfseries END OF EXAMINATION}
\end{center}

\end{document}
